\documentclass[11pt]{article}
\usepackage[margin=1in]{geometry}
\usepackage{amsmath}
\usepackage{amssymb}
\usepackage{amsthm}
\usepackage{algorithm}
\usepackage{algpseudocode}
\usepackage{booktabs}
\usepackage{graphicx}
\usepackage{url}
\usepackage[hidelinks]{hyperref}

\newtheorem{theorem}{Theorem}
\newtheorem{lemma}{Lemma}
\newtheorem{corollary}{Corollary}
\theoremstyle{definition}
\newtheorem{definition}{Definition}

\newcommand{\F}{\mathbb{F}}
\newcommand{\R}{\mathbf{R}}
\newcommand{\Rk}[2]{\mathbf{R}_{#1}(#2)}
\newcommand{\Cyct}[1]{\mathrm{C}_{#1}}
\newcommand{\Fullt}[1]{\mathrm{P}_{#1}}
\newcommand{\Trunct}[1]{\mathrm{T}_{#1}}
\newcommand{\Negat}[1]{\mathrm{C}^{-}_{#1}}
\newcommand{\GL}{\mathrm{GL}}
\newcommand{\PGL}{\mathrm{PGL}}
\newcommand{\Gal}{\mathrm{Gal}}
\newcommand{\Aut}{\mathrm{Aut}}
\newcommand{\Kappa}{\mathrm{K}}
\newcommand{\rank}{\operatorname{rank}}
\newcommand{\RREF}{\operatorname{RREF}}
\newcommand{\mmt}[3]{\langle #1,#2,#3\rangle}

\algnewcommand\algorithmicparallelfor{\textbf{parallel for}}
\algdef{SE}[PARFOR]{ParallelFor}{EndParallelFor}[1]%
{\algorithmicparallelfor\ #1\ \algorithmicdo}{\algorithmicend\ \algorithmicparallelfor}

\setlength{\parindent}{0pt}
\setlength{\parskip}{0.6\baselineskip}

\bibliographystyle{plain}

\title{Automated Lower Bounds for Bilinear Complexity\\ over Finite Fields}
\author{Chengu Wang\\ \texttt{wangchengu@gmail.com}}
\date{}

\begin{document}
\maketitle

\begin{abstract}
  We present a general, automated framework for proving lower bounds on the
  bilinear complexity (tensor rank) of multiplication problems over a finite field
  $\mathbb{F}_q$. The framework is parameterized only by the multiplication tensor and by a
  group of rank-preserving symmetries acting on one argument: it classifies
  the orbits of constraint subspaces under that group, runs a dynamic program over
  the orbits combining four lower-bound techniques, and emits a proof certificate
  that a verifier rechecks, typically faster than the search.

  Instantiating the framework for matrix multiplication, we improve the lower bounds
  for four small formats over $\mathbb{F}_2$, most notably showing that
  the bilinear complexity of multiplying two $3 \times 3$ matrices over $\mathbb{F}_2$ is at
  least $20$, raising the bound of $19$ that had stood since Bl\"aser (2003).
  Instantiating it for polynomial multiplication --- full products,
  cyclic convolution, and the truncated (modulo $x^N$) and negacyclic (modulo
  $x^N+1$) products --- we obtain eighteen new lower bounds over $\mathbb{F}_2$ and $\mathbb{F}_3$.
  Every bound is backed by a machine-checkable certificate.
\end{abstract}

\section{Introduction}\label{sec:intro}

A \emph{bilinear} algorithm for a bilinear map $\beta\colon U\times V\to W$ over a
field $\F$ computes a fixed set of products of a linear form in the $U$-variables
with a linear form in the $V$-variables, and then reads off each coordinate of
the output as a linear combination of those products. The minimum number of such
products is the \emph{bilinear complexity} of $\beta$, and it equals the rank of
the order-$3$ tensor naturally associated with $\beta$. Two families of bilinear
maps have driven the subject: matrix multiplication and polynomial
multiplication. For matrix multiplication the central quantity is
$\R(\mmt{l}{m}{n})$, the rank of the tensor for multiplying an $l\times m$ by an
$m\times n$ matrix; for polynomial multiplication it is the rank of the tensor
for multiplying two polynomials.

Over small finite fields these ranks are notoriously hard to pin down, and even
for tiny formats the best known lower and upper bounds often disagree. This paper
develops one automated method that attacks both families, and uses it to push
several of these small cases.

\subsection{Lower bounds for matrix multiplication}

Strassen~\cite{strassen1969gaussian} showed $\R(\mmt{2}{2}{2})\le 7$ over $\mathbb{Z}$, and
Winograd~\cite{winograd1971multiplication} proved the matching lower bound over $\mathbb{Q}$.
Then, Hopcroft and Kerr~\cite{hopcroft1971minimizing} established several lower
bounds over $\F_2$, including $\Rk{\F_2}{\mmt{2}{3}{3}}\ge 15$, and stated
$\Rk{\F_2}{\mmt{2}{3}{4}}\ge 19$ without a full proof.
Bshouty~\cite{bshouty1989lower} proved $\Rk{\F_2}{\mmt{n}{n}{n}}\ge \tfrac{5}{2}n^2-o(n^2)$.
Bl\"aser established the influential bounds
$\R(\mmt{l}{m}{n})\ge lm+mn+l-m+n-3$ over algebraically closed
fields~\cite{blaser1999lower}, $\R(\mmt{n}{n}{n})\ge \tfrac{5}{2}n^2-3n$ over
arbitrary fields~\cite{blaser1999fivehalf}, and finally
$\R(\mmt{n}{m}{n})\ge 2mn+2n-m-2$ for $m\ge n\ge 3$~\cite{blaser2003complexity} over arbitrary fields,
which implies $\R(\mmt{3}{3}{3})\ge 19$. On the upper side
Laderman~\cite{laderman1976noncommutative} multiplied $3\times 3$ matrices with
$23$ products over $\mathbb{Z}$. Thus the rank of $\mmt{3}{3}{3}$ has remained trapped in $[19,23]$,
making it the prototypical open small format.

\subsection{Lower bounds for polynomial multiplication}

The bilinear complexity of polynomial multiplication over finite fields has an
equally long history, and unlike matrix multiplication it is governed by the
arithmetic of $\F_q$. The \emph{full} product is the kernel of fast integer and
polynomial multiplication, where recursive schemes reduce a large product to many
small base cases whose cost compounds. For two degree-$(N-1)$ polynomials, Fiduccia
and Zalcstein~\cite{fiduccia1977algebras} gave the lower
bound $2N-1$, met by Toom--Cook evaluation--interpolation whenever the field has
enough points ($2N-1\le q+1$)~\cite{toom1963complexity}. When the field is too small for plain
interpolation the best known upper bounds come instead from Karatsuba-like
formulae, notably those of Montgomery~\cite{montgomery2005five}, and from
CRT/evaluation--interpolation over higher-degree places~\cite{winograd1977some}.
Outside that range the
exact value climbs: Kaminski and Bshouty~\cite{kaminski1989multiplicative}
determined it precisely in the next range, and for still larger $N$ the strongest
lower bounds come from embedding a bilinear algorithm into a linear code and
invoking the Griesmer bound~\cite{griesmer1960bound} or the tables of
Grassl~\cite{grassl2007codetables}, via the code constructions of
Lempel--Seroussi--Winograd~\cite{lempel1983complexity} and
Chudnovsky--Chudnovsky~\cite{chudnovsky1988algebraic}; small cases were settled by
the exhaustive search of Barbulescu, Detrey, Estibals, and
Zimmermann~\cite{barbulescu2012finding}.

A second strand is multiplication in a polynomial quotient ring
$\F_q[x]/(x^N-\gamma)$, which specializes to three classical families:
\emph{cyclic} convolution ($\gamma=1$, modulo $x^N-1$), multiplication in the group
algebra of a cyclic group, underlying the discrete Fourier transform, digital signal
processing, and cyclic codes; the \emph{truncated} (or short) product ($\gamma=0$,
modulo $x^N$), multiplication in the local algebra $\F_q[x]/x^N$ that is the
recurring base case of Karatsuba- and Toom-style schemes; and the \emph{negacyclic}
product ($\gamma=-1$, modulo $x^N+1$), a twisted group algebra underlying skew-cyclic
codes and lattice-based cryptography. The lower bounds for all three are governed by
the factorization of $x^N-\gamma$: Winograd~\cite{winograd1977some} proved that
multiplication modulo a polynomial with $k$ distinct irreducible factors needs at
least $2N-k$ products, and Bl\"aser's characterization of algebras of minimal
bilinear complexity~\cite{blaser2005complete}, combined with the minimal-rank theory
of Averbuch, Galil, and
Winograd~\cite{averbuch1988classification,averbuch1991classification} (whose Part~II
treats exactly the truncated algebra $\F_q[x]/x^N$), strengthens these on the
local-algebra factors. Upper bounds come from the convolution constructions of Wagh
and Morgera~\cite{wagh1983structured} and Morgera~\cite{morgera1990multiplicative},
refined by Cenk and \"Ozbudak~\cite{cenk2010multiplication,cenk2011multiplication}.
Over characteristic $2$ the negacyclic product coincides with the cyclic one, since
$x^N+1=x^N-1$.

\subsection{Our contributions}

We give a single framework, described abstractly in
Section~\ref{sec:framework} and instantiated per problem in
Section~\ref{sec:symmetries}, that proves bilinear-complexity lower bounds for any
multiplication tensor equipped with a group of first-argument symmetries. It
classifies the orbits of constraint subspaces (Section~\ref{sec:orbits}), runs a
dimension-sweeping dynamic program with four lower-bound techniques
(Section~\ref{sec:techniques}), and produces certificates that an independent
verifier rechecks deterministically (Section~\ref{sec:certificates}).

Our main results concern matrix multiplication, where the framework improves the
best known lower bound for four small formats over $\F_2$.

\begin{theorem}[Matrix multiplication]\label{thm:matrix}
  \[
    \Rk{\F_2}{\mmt{2}{3}{4}}\ge 19,\quad
    \Rk{\F_2}{\mmt{3}{3}{3}}\ge 20,\quad
    \Rk{\F_2}{\mmt{3}{3}{4}}\ge 25,\quad
    \Rk{\F_2}{\mmt{3}{4}{4}}\ge 29.
  \]
\end{theorem}

The headline is $\Rk{\F_2}{\mmt{3}{3}{3}}\ge 20$, which raises the value $19$ that had
stood since Bl\"aser (2003)~\cite{blaser2003complexity}; the search finds this proof
in about $40$ minutes on a laptop and the certificate verifies in seconds. The same
certificate has since been re-verified independently by
Beuchert~\cite{beuchert_2026_20752782}, using a checker written from scratch in a
different programming language. The bound
$\Rk{\F_2}{\mmt{2}{3}{4}}\ge 19$ completes a bound stated without proof by Hopcroft and
Kerr~\cite{hopcroft1971minimizing}. Table~\ref{tab:matrix} places all four in
context.

\begin{table}[h]
  \centering
  \setlength{\tabcolsep}{14pt}
  \begin{tabular}{llll}
    \toprule
    format          & prev LB                                                          & our LB       & prev UB                                                      \\
    \midrule
    $\mmt{2}{2}{2}$ & $7$~\cite{winograd1971multiplication}                            & $7$          & $7$~\cite{strassen1969gaussian}                              \\
    $\mmt{2}{2}{3}$ & $11$~\cite{hopcroft1971minimizing}                               & $11$         & $11$ ($\langle 2, 2, 2:7\rangle + \langle 2, 2, 1:4\rangle$) \\
    $\mmt{2}{2}{4}$ & $14$~\cite{hopcroft1971minimizing}                               & $14$         & $14$ ($\langle 2, 2, 2:7\rangle + \langle 2, 2, 2:7\rangle$) \\
    $\mmt{2}{3}{3}$ & $15$~\cite{hopcroft1971minimizing}                               & $15$         & $15$~\cite{hopcroft1971minimizing}                           \\
    $\mmt{2}{3}{4}$ & $18$~\cite{blaser1999lower}, $19$?~\cite{hopcroft1971minimizing} & $\boxed{19}$ & $20$~\cite{hopcroft1971minimizing}                           \\
    $\mmt{3}{3}{3}$ & $19$~\cite{blaser2003complexity}                                 & $\boxed{20}$ & $23$~\cite{laderman1976noncommutative}                       \\
    $\mmt{3}{3}{4}$ & $24$~\cite{blaser2003complexity}                                 & $\boxed{25}$ & $29$~\cite{smirnov2013bilinear}                              \\
    $\mmt{3}{4}{4}$ & $28$~\cite{blaser1999lower}                                      & $\boxed{29}$ & $38$~\cite{smirnov2013bilinear}                              \\
    \bottomrule
  \end{tabular}
  \caption{Matrix multiplication tensor rank over $\F_2$; boxed entries are new.}
  \label{tab:matrix}
\end{table}

For polynomial multiplication the same framework yields eighteen new lower bounds,
over $\F_2$ and $\F_3$ --- for the full product and for the
three quotient families $\F_q[x]/(x^N-\gamma)$.

\begin{theorem}[Polynomial multiplication]\label{thm:poly}
  Write $\Rk{\F_q}{\Fullt{N}}$, $\Rk{\F_q}{\Cyct{N}}$, $\Rk{\F_q}{\Trunct{N}}$, and
  $\Rk{\F_q}{\Negat{N}}$ for the bilinear complexity over $\F_q$ of, respectively, the
  full product of two degree-$(N-1)$ polynomials, cyclic convolution (modulo $x^N-1$),
  the truncated product (modulo $x^N$), and the negacyclic product (modulo $x^N+1$),
  each of length $N$. Then
  \[
    \begin{aligned}
      \text{full:}       & \quad \Rk{\F_2}{\Fullt{6}}\ge 16,\ \ \Rk{\F_2}{\Fullt{7}}\ge 19,\ \ \Rk{\F_2}{\Fullt{8}}\ge 21,\ \ \Rk{\F_3}{\Fullt{6}}\ge 14;                                      \\
      \text{cyclic:}     & \quad \Rk{\F_2}{\Cyct{7}}\ge 13,\ \ \Rk{\F_2}{\Cyct{8}}\ge 19,\ \ \Rk{\F_2}{\Cyct{10}}\ge 22,\ \ \Rk{\F_3}{\Cyct{9}}\ge 19;                                         \\
      \text{truncated:}  & \quad \Rk{\F_2}{\Trunct{6}}\ge 13,\ \ \Rk{\F_2}{\Trunct{7}}\ge 16,\ \ \Rk{\F_2}{\Trunct{8}}\ge 19,\ \ \Rk{\F_2}{\Trunct{9}}\ge 21,\ \ \Rk{\F_2}{\Trunct{10}}\ge 24, \\
                         & \quad \Rk{\F_2}{\Trunct{11}}\ge 27,\ \ \Rk{\F_3}{\Trunct{8}}\ge 17,\ \ \Rk{\F_3}{\Trunct{9}}\ge 19,\ \ \Rk{\F_3}{\Trunct{10}}\ge 21;                                \\
      \text{negacyclic:} & \quad \Rk{\F_3}{\Negat{9}}\ge 19.
    \end{aligned}
  \]
\end{theorem}

Tables~\ref{tab:full},~\ref{tab:cyc},~\ref{tab:trunc}, and~\ref{tab:nega} place these
against the previous best lower bound and the best known upper bound; the boxed
entries are the improvements. In one case (cyclic length $7$ over $\F_2$) the
new lower bound meets the best known upper bound, settling
$\Rk{\F_2}{\Cyct{7}}=13$. Over $\F_2$ the truncated product also yields a new upper
bound, $\Rk{\F_2}{\Trunct{9}}\le 26$, found by the flip-graph search~\cite{kauers2023flip}.
Over characteristic $2$ the negacyclic product coincides with the
cyclic one ($x^N+1=x^N-1$), so we skip $\F_2$ in Table~\ref{tab:nega}.

\begin{table}[h]
  \centering
  \setlength{\tabcolsep}{6pt}
  \begin{tabular}{lllllll}
    \toprule
        & \multicolumn{3}{c}{over $\F_2$}                       & \multicolumn{3}{c}{over $\F_3$}                                                                                                                                                                \\
    \cmidrule(lr){2-4}\cmidrule(lr){5-7}
    $N$ & prev LB                                               & our LB                          & prev UB                                         & prev LB                                            & our LB       & prev UB                                \\
    \midrule
    $1$ & $1$                                                   & $1$                             & $1$                                             & $1$                                                & $1$          & $1$                                    \\
    $2$ & $3$~\cite{fiduccia1977algebras}                       & $3$                             & $3$~\cite{toom1963complexity}                   & $3$~\cite{fiduccia1977algebras}                    & $3$          & $3$~\cite{toom1963complexity}          \\
    $3$ & $6$~\cite{lempel1983complexity,griesmer1960bound}     & $6$                             & $6$~\cite{kaminski1989multiplicative}           & $6$~\cite{kaminski1989multiplicative}              & $6$          & $6$~\cite{kaminski1989multiplicative}  \\
    $4$ & $9$~\cite{kaminski1989multiplicative}                 & $9$                             & $9$~\cite{kaminski1989multiplicative}           & $9$~\cite{kaminski1989multiplicative}              & $9$          & $9$~\cite{kaminski1989multiplicative}  \\
    $5$ & $13$~\cite{lempel1983complexity,grassl2007codetables} & $13$                            & $13$~\cite{montgomery2005five}                  & $12$~\cite{kaminski1989multiplicative}             & $12$         & $12$~\cite{kaminski1989multiplicative} \\
    $6$ & $15$~\cite{lempel1983complexity,grassl2007codetables} & $\boxed{16}$                    & $17$~\cite{montgomery2005five}                  & $12$~\cite{lempel1983complexity,griesmer1960bound} & $\boxed{14}$ & $15$~\cite{winograd1977some}           \\
    $7$ & $18$~\cite{lempel1983complexity,grassl2007codetables} & $\boxed{19}$                    & $22$~\cite{montgomery2005five}                  &                                                    &              &                                        \\
    $8$ & $20$~\cite{lempel1983complexity,grassl2007codetables} & $\boxed{21}$                    & $26$~\cite{winograd1977some,fanhasan2015survey} &                                                    &              &                                        \\
    \bottomrule
  \end{tabular}
  \caption{Full product of two degree-$(N-1)$ polynomials; boxed entries are new.}
  \label{tab:full}
\end{table}

\begin{table}[h]
  \centering
  \setlength{\tabcolsep}{6pt}
  \begin{tabular}{lllllll}
    \toprule
         & \multicolumn{3}{c}{over $\F_2$} & \multicolumn{3}{c}{over $\F_3$}                                                                                                                                                                               \\
    \cmidrule(lr){2-4}\cmidrule(lr){5-7}
    $N$  & prev LB                         & our LB                          & prev UB                                                      & prev LB                        & our LB       & prev UB                                                      \\
    \midrule
    $1$  & $1$                             & $1$                             & $1$                                                          & $1$                            & $1$          & $1$                                                          \\
    $2$  & $3$~\cite{winograd1977some}     & $3$                             & $3$~\cite{morgera1990multiplicative}                         & $2$~\cite{winograd1977some}    & $2$          & $2$~\cite{wagh1983structured}                                \\
    $3$  & $4$~\cite{winograd1977some}     & $4$                             & $4$~\cite{wagh1983structured}                                & $5$~\cite{winograd1977some}    & $5$          & $5$~\cite{morgera1990multiplicative}                         \\
    $4$  & $8$~\cite{blaser2005complete}   & $8$                             & $8$~\cite{morgera1990multiplicative,cenk2011multiplication}  & $5$~\cite{winograd1977some}    & $5$          & $5$~\cite{wagh1983structured}                                \\
    $5$  & $10$~\cite{blaser2005complete}  & $9$                             & $10$~\cite{wagh1983structured}                               & $10$~\cite{blaser2005complete} & $9$          & $10$~\cite{wagh1983structured}                               \\
    $6$  & $11$~\cite{blaser2005complete}  & $11$                            & $12$~\cite{morgera1990multiplicative}                        & $10$~\cite{winograd1977some}   & $10$         & $10$~\cite{barbulescu2012finding}                            \\
    $7$  & $12$~\cite{blaser2005complete}  & $\boxed{13}$                    & $13$~\cite{wagh1983structured}                               & $13$~\cite{blaser2005complete} & $13$         & $16$~\cite{morgera1990multiplicative,cenk2010multiplication} \\
    $8$  & $16$~\cite{blaser2005complete}  & $\boxed{19}$                    & $22$~\cite{morgera1990multiplicative,cenk2011multiplication} & $11$~\cite{winograd1977some}   & $11$         & $11$~\cite{wagh1983structured}                               \\
    $9$  & $19$~\cite{blaser2005complete}  & $18$                            & $19$~\cite{morgera1990multiplicative}                        & $18$~\cite{blaser2005complete} & $\boxed{19}$ & $27$~\cite{morgera1990multiplicative,cenk2011multiplication} \\
    $10$ & $19$~\cite{blaser2005complete}  & $\boxed{22}$                    & $29$~\cite{morgera1990multiplicative}                        &                                &              &                                                              \\
    \bottomrule
  \end{tabular}
  \caption{Cyclic convolution modulo $x^N-1$; boxed entries are new.}
  \label{tab:cyc}
\end{table}

\begin{table}[h]
  \centering
  \setlength{\tabcolsep}{6pt}
  \begin{tabular}{lllllll}
    \toprule
         & \multicolumn{3}{c}{over $\F_2$}                           & \multicolumn{3}{c}{over $\F_3$}                                                                                                                                                             \\
    \cmidrule(lr){2-4}\cmidrule(lr){5-7}
    $N$  & prev LB                                                   & our LB                          & prev UB                                   & prev LB                                                   & our LB       & prev UB                            \\
    \midrule
    $1$  & $1$                                                       & $1$                             & $1$                                       & $1$                                                       & $1$          & $1$                                \\
    $2$  & $3$~\cite{averbuch1991classification,blaser2005complete}  & $3$                             & $3$~\cite{toom1963complexity}             & $3$~\cite{averbuch1991classification,blaser2005complete}  & $3$          & $3$~\cite{toom1963complexity}      \\
    $3$  & $5$~\cite{averbuch1991classification,blaser2005complete}  & $5$                             & $5$~\cite{cenk2011multiplication}         & $5$~\cite{averbuch1991classification,blaser2005complete}  & $5$          & $5$~\cite{cenk2011multiplication}  \\
    $4$  & $8$~\cite{averbuch1991classification,blaser2005complete}  & $8$                             & $8$~\cite{cenk2011multiplication}         & $8$~\cite{averbuch1991classification,blaser2005complete}  & $7$          & $8$~\cite{cenk2011multiplication}  \\
    $5$  & $10$~\cite{averbuch1991classification,blaser2005complete} & $10$                            & $11$~\cite{cenk2011multiplication}        & $10$~\cite{averbuch1991classification,blaser2005complete} & $10$         & $10$~\cite{barbulescu2012finding}  \\
    $6$  & $12$~\cite{averbuch1991classification,blaser2005complete} & $\boxed{13}$                    & $14$~\cite{cenk2011multiplication}        & $12$~\cite{averbuch1991classification,blaser2005complete} & $12$         & $14$~\cite{cenk2011multiplication} \\
    $7$  & $14$~\cite{averbuch1991classification,blaser2005complete} & $\boxed{16}$                    & $18$~\cite{cenk2011multiplication}        & $14$~\cite{averbuch1991classification,blaser2005complete} & $14$         & $18$~\cite{cenk2011multiplication} \\
    $8$  & $16$~\cite{averbuch1991classification,blaser2005complete} & $\boxed{19}$                    & $22$~\cite{cenk2011multiplication}        & $16$~\cite{averbuch1991classification,blaser2005complete} & $\boxed{17}$ & $22$~\cite{cenk2011multiplication} \\
    $9$  & $18$~\cite{averbuch1991classification,blaser2005complete} & $\boxed{21}$                    & $27$~\cite{cenk2011multiplication}$^\ast$ & $18$~\cite{averbuch1991classification,blaser2005complete} & $\boxed{19}$ & $27$~\cite{winograd1977some}       \\
    $10$ & $20$~\cite{averbuch1991classification,blaser2005complete} & $\boxed{24}$                    & $31$~\cite{cenk2011multiplication}        & $20$~\cite{averbuch1991classification,blaser2005complete} & $\boxed{21}$ & $31$~\cite{winograd1977some}       \\
    $11$ & $22$~\cite{averbuch1991classification,blaser2005complete} & $\boxed{27}$                    & $36$~\cite{cenk2011multiplication}        &                                                           &              &                                    \\
    \bottomrule
  \end{tabular}

  \smallskip
  {\footnotesize\raggedright
    $^\ast$ We obtain a new upper bound of $26$ by flip-graph search~\cite{kauers2023flip}.
    \scriptsize
    $a0\otimes (b0+b2+b4+b6+b8)\otimes c8 + a1\otimes (b1+b3+b4+b5+b6+b7)\otimes (c7+c8) + (a0+a1)\otimes (b0+b1)\otimes (c1+c2+c4) + (a0+a1)\otimes (b0+b1+b3+b7)\otimes c7 + a3\otimes (b0+b4)\otimes (c4+c7) + (a0+a1+a4)\otimes (b1+b4)\otimes (c4+c5+c7) + (a0+a1+a4)\otimes (b0+b3+b4)\otimes (c5+c7+c8) + (a2+a4)\otimes (b2+b3)\otimes (c5+c6) + (a3+a4)\otimes (b0+b1+b2+b4)\otimes (c4+c6) + (a2+a3+a4)\otimes (b0+b2+b3+b4)\otimes (c4+c5+c6+c7) + (a1+a2+a3+a4)\otimes (b0+b3+b4)\otimes (c4+c5+c7) + (a4+a5)\otimes (b0+b1+b3+b5)\otimes (c3+c5) + (a0+a4+a5)\otimes (b1+b3+b5)\otimes (c3+c5+c8) + (a0+a3+a4+a5)\otimes (b0+b1+b5)\otimes (c3+c6) + (a0+a1+a3+a4+a5)\otimes (b0+b5)\otimes (c3+c6+c7+c8) + (a0+a2+a3+a4+a5)\otimes (b0+b1+b2+b5)\otimes (c3+c7) + (a1+a2+a3+a4+a5)\otimes (b2+b5)\otimes c3 + a6\otimes (b0+b2+b6)\otimes (c2+c3+c7) + (a0+a6)\otimes (b2+b6)\otimes (c2+c3+c7+c8) + (a1+a2+a6)\otimes (b1+b6)\otimes (c2+c3+c6+c7+c8) + (a0+a1+a2+a6)\otimes (b0+b1+b6)\otimes (c2+c3+c6+c7) + (a0+a1+a3+a4+a7)\otimes (b0+b1)\otimes c7 + (a0+a2+a4+a5+a6+a7)\otimes b1\otimes (c1+c5+c7+c8) + (a0+a1+a2+a4+a5+a6+a7)\otimes b1\otimes (c1+c5) + (a3+a5+a8)\otimes b0\otimes (c0+c1+c4+c5+c6) + (a0+a3+a5+a8)\otimes b0\otimes (c0+c1+c4+c5+c6+c8)$
    \par}
  \caption{Truncated product modulo $x^N$; boxed entries are new lower bounds.}
  \label{tab:trunc}
\end{table}

\begin{table}[h]
  \centering
  \setlength{\tabcolsep}{12pt}
  \begin{tabular}{llll}
    \toprule
        & \multicolumn{3}{c}{over $\F_3$}                                                          \\
    \cmidrule(lr){2-4}
    $N$ & prev LB                         & our LB       & prev UB                                 \\
    \midrule
    $1$ & $1$                             & $1$          & $1$                                     \\
    $2$ & $3$~\cite{winograd1977some}     & $3$          & $3$~\cite{toom1963complexity}           \\
    $3$ & $5$~\cite{winograd1977some}     & $5$          & $5$~\cite{morgera1990multiplicative}    \\
    $4$ & $6$~\cite{winograd1977some}     & $6$          & $6$~\cite{degroote1983characterization} \\
    $5$ & $10$~\cite{blaser2005complete}  & $9$          & $10$~\cite{cenk2010multiplication}      \\
    $6$ & $12$~\cite{blaser2005complete}  & $12$         & $15$~\cite{morgera1990multiplicative}   \\
    $7$ & $13$~\cite{blaser2005complete}  & $13$         & $16$~\cite{cenk2010multiplication}      \\
    $8$ & $16$~\cite{blaser2005complete}  & $16$         & $18$~\cite{cenk2010multiplication}      \\
    $9$ & $18$~\cite{blaser2005complete}  & $\boxed{19}$ & $27$~\cite{cenk2011multiplication}      \\
    \bottomrule
  \end{tabular}
  \caption{Negacyclic convolution modulo $x^N+1$; the boxed entry is
    new. Over $\F_2$, $x^N+1=x^N-1$, so the negacyclic product coincides with cyclic
    convolution (Table~\ref{tab:cyc}).}
  \label{tab:nega}
\end{table}

The rest of the paper is organized as follows.
Section~\ref{sec:framework} sets up tensors, constraints, and symmetries. Section~\ref{sec:symmetries} instantiates the symmetries for
matrix and polynomial multiplication. Section~\ref{sec:orbits} describes orbit
enumeration.
Section~\ref{sec:techniques} formalizes the four lower-bound techniques.
Section~\ref{sec:certificates} details the certificates and the verifier, and
Section~\ref{sec:conclusion} concludes. Appendix~\ref{app:results} reports the
search and verification details for each new bound.

\section{The framework}\label{sec:framework}

\subsection{Tensors, rank, and flattening}

Fix a finite field $\F_q$ with $q=p^m$. A multiplication problem is given by an
order-$3$ tensor
\[
  T = \sum_{i,j,k} t_{ijk}\, a_i \otimes b_j \otimes c_k
  \;\in\; A \otimes B \otimes C,
\]
where $A=\F_q^{N_A}$, $B=\F_q^{N_B}$, $C=\F_q^{N_C}$ have bases $a_i,b_j,c_k$ and
the structure constants $t_{ijk}\in\F_q$ encode the bilinear map. The
\emph{rank} $\R(T)$ is the least $r$ such that
$T=\sum_{\lambda=1}^{r} u_\lambda\otimes v_\lambda\otimes w_\lambda$ with
$u_\lambda\in A$, $v_\lambda\in B$, $w_\lambda\in C$; this is exactly the bilinear
complexity of the underlying map. We write $\Rk{\F_q}{T}$ when emphasizing the field $\F_q$,
and $\R(T)$ when it is clear from context.

The one universal lower bound is \emph{flattening}: grouping two of the three
tensor factors into one turns $T$ into a matrix whose rank lower-bounds $\R(T)$.
Concretely, $\R(T)\ge \rank M_A(T)$ where $M_A(T)$ is the $N_A\times (N_B N_C)$
matrix with rows indexed by $A$ and columns by $B\otimes C$; symmetrically for the
other two groupings, and we take the maximum.

\subsection{Constraints}

We attack $\R(T)$ by constraining the first argument $A$ to a subspace. A
\emph{constraint} is a linear functional on $A$, i.e.\ a vector in the dual
$A^* = \F_q^{N_A}$; a set of constraints cuts $A$ down to the subspace on which
all of them vanish, and substituting that subspace into $T$ yields a smaller
tensor $T_S$ with $\R(T)\ge \R(T_S)$ for any subspace $S$ obtained by adding
constraints. A set of $d$ independent constraints spans a $d$-dimensional
subspace of $A^*$, and two constraint sets that span the same subspace define the
same $T_S$. We therefore represent a constraint subspace by its
\emph{reduced row echelon form} (RREF), which is a unique key for the subspace.

\subsection{Symmetries and the query/store factorization}

The leverage in our method comes from symmetries that act on $A$ and preserve
$\R(T)$. Abstractly, let $G$ be a group acting on $A^*$ by semilinear
bijections: each $g$ is additive with $g(\lambda u)=\phi_g(\lambda)\,g(u)$ for
some field automorphism $\phi_g\in\Gal(\F_q/\F_p)$, so both the $\F_q$-linear
sandwich and multiplication maps ($\phi_g=\mathrm{id}$) and the Galois symmetries
($\phi_g$ a Frobenius $a\mapsto a^p$) are allowed. We require that for every
$g\in G$ there is a rank-preserving symmetry of $T$ whose action on the first
factor is $g$. Since a semilinear bijection sends $\F_q$-subspaces to
$\F_q$-subspaces, $G$ permutes constraint subspaces and maps each $T_S$ to an
isomorphic tensor, so $\R(T_S)$ is constant on each $G$-orbit of subspaces. Our
dynamic program needs only one representative per orbit.

The implementation never materializes $G$ as a list. Instead it factors the group $G$
into two subsets of $G$, a \emph{query} set $\Kappa$ and a \emph{store} set
$\Sigma$, such that the product set
$\Kappa^{-1}\cdot\Sigma$ covers $G$ so that every orbit member is reached.
Neither $\Kappa$ nor $\Sigma$ needs to be a subgroup or be closed under inversion.
This is the meet-in-the-middle, or
square-root, trick. To test whether a query subspace $g$ lies in the orbit of a
stored representative $g'$, we precompute and hash the RREF of
$\sigma(g')$ for every $\sigma\in\Sigma$ (the \emph{store images}), and
at query time probe the hash with the RREF of $\kappa(g)$ for every
$\kappa\in\Kappa$. A hit
$\kappa(g)\equiv \sigma(g')$ certifies $g$ and $g'$ as
$G$-equivalent and yields the witness pair $(\kappa,\sigma)$, from which the canonical form
is recovered as
\begin{equation}\label{eq:witness}
  g' \;\equiv\; \sigma^{-1}\bigl(\kappa(g)\bigr).
\end{equation}
Storage grows by a factor $|\Sigma|$ and a query costs $|\Kappa|$
probes, with $|\Kappa|\cdot|\Sigma|\gtrsim |G|$; choosing
$|\Kappa|\approx|\Sigma|\approx\sqrt{|G|}$ trades square-root space for
square-root time.

\section{Symmetries of the problem families}\label{sec:symmetries}

The framework is fixed once we name, for each problem, the tensor $T$ and the
factorization $(\Kappa,\Sigma)$ of its first-argument symmetry group.

\subsection{Matrix multiplication}

The tensor for multiplying an $l\times m$ matrix $X$ by an $m\times n$ matrix $Y$, with product $Z^{\!\top}$, is
\[
  \mmt{l}{m}{n} = \sum_{i=0}^{l-1}\sum_{j=0}^{m-1}\sum_{k=0}^{n-1}
  x_{ij}\otimes y_{jk}\otimes z_{ki},
\]
so $A=\F_q^{l\times m}$. For invertible $P\in\GL_l$, $Q\in\GL_m$, $R\in\GL_n$ the
substitution $X\mapsto PXQ^{-1}$, $Y\mapsto QYR^{-1}$, $Z\mapsto RZP^{-1}$ preserves
the tensor (the \emph{sandwich} symmetry). The rank is also invariant under
cyclically permuting the three factors, $\R(\mmt{l}{m}{n})=\R(\mmt{m}{n}{l})$ (the
\emph{cyclic} symmetry); for a square format $l=m=n$ this combines with matrix
transposition into an order-two symmetry
$X\mapsto X^{\!\top}$, $Y\mapsto Z^{\!\top}$, $Z\mapsto Y^{\!\top}$ (the \emph{transpose} symmetry). Projected onto the
first factor, the symmetry group acting on the matrix space $A=\F_q^{l\times m}$ is
\[
  G \;=\; (\GL_l(\F_q)\times \GL_m(\F_q)) \rtimes C_2,
\]
with $(P,Q)$ acting by $X\mapsto PXQ^{-1}$ and the $C_2$ by transposition $X\mapsto X^{\!\top}$ (present
only when $l=m=n$). The meet-in-the-middle splits $G$ by side:
the store set is the right multiplications $\Sigma=\{X\mapsto XR^{-1}:R\in\GL_m(\F_q)\}$ and
the query set is the left multiplications $\Kappa=\{X\mapsto LX:L\in\GL_l(\F_q)\}$,
extended by the transpose when $l=m=n$. Scalar matrices act trivially on subspaces,
so each side is taken modulo the center, and both then have order about $\sqrt{|G|}$.
Equivalence of constraint subspaces under the subgroup without
transpose is exactly the tensor-isomorphism problem on $l\times m\times d$ tensors,
which sits above graph and code equivalence in the isomorphism
hierarchy~\cite{grochow2023complexity,grochow2025complexity}.

\subsection{Full polynomial multiplication}

Multiplying two polynomials of degree $<N$ in $\F_q[x]$ gives
\[
  \Fullt{N} = \sum_{i,j=0}^{N-1} a_i\otimes b_j\otimes c_{i+j},
  \qquad N_A=N_B=N,\ N_C=2N-1,
\]
the ordinary convolution. Its first-argument symmetries are the action on
degree-$(N-1)$ binary forms of the projective semilinear group
\[
  G_{\mathrm{full}} = \mathrm{P\Gamma L}_2(\F_q)
  = \PGL_2(\F_q)\rtimes \Gal(\F_q/\F_p),
\]
generated by the substitutions $x\mapsto x+c$, the reversal $x\mapsto 1/x$ (which
swaps $a_i\leftrightarrow a_{N-1-i}$), the scalings $x\mapsto gx$, and the
Frobenius $a\mapsto a^p$. We factor it as
$\Sigma=\PGL_2(\F_q)$ (the store side, of size $q(q^2-1)$) and
$\Kappa=\Gal(\F_q/\F_p)$ (the query side, of size $m$).

\subsection{Polynomial multiplication modulo $x^N-\gamma$}

Fix a scalar $\gamma\in\F_q$ and let $R=\F_q[x]/(x^N-\gamma)$. Multiplication in $R$
reduces the convolution of $a$ and $b$ by the rule $x^N\equiv\gamma$,
\[
  \sum_{i,j=0}^{N-1} a_i\otimes b_j\otimes w_{ij},\qquad
  w_{ij}=\begin{cases} c_{i+j},           & i+j<N,    \\[2pt]
                \gamma\,c_{i+j-N}, & i+j\ge N,\end{cases}
  \qquad N_A=N_B=N_C=N.
\]
Three families specialize this tensor: \emph{cyclic} convolution $\Cyct{N}$ for
$\gamma=1$, the \emph{truncated} (or short) product $\Trunct{N}$ for $\gamma=0$ ---
the local algebra $\F_q[x]/(x^N)$, in which the overflow terms vanish and $x$ is
nilpotent --- and the \emph{negacyclic} product $\Negat{N}$ for $\gamma=-1$. Over
characteristic $2$, $-1=1$, so $\Negat{N}=\Cyct{N}$.

The same three kinds of symmetry act on $A=R$ for every $\gamma$: multiplication by a
unit $p\in R^*$ (the shifts $x^r$ are the special case $p=x^r$, valid for
$\gamma\ne 0$), the $\F_q$-algebra automorphisms $\Aut(R)$, and the Galois group of
$\F_q/\F_p$. The base-field scalars $\F_q^*$ act trivially on subspaces, so the group
acting on constraint subspaces is
\[
  G = \bigl(R^*\rtimes \Aut(R) \rtimes \Gal(\F_q/\F_p)\bigr)/\F_q^*,
\]
with store side $\Sigma=R^*/\F_q^*$ and query side
$\Kappa=\Aut(R)\rtimes\Gal(\F_q/\F_p)$ --- the identical factorization in all three
cases. Only the ring structure, and hence the sizes of the two sides, depends on
$\gamma$. For $\gamma=\pm1$ the Chinese remainder theorem splits $R$ into field
factors --- for $\gamma=1$ indexed by the $q$-cyclotomic cosets of $\mathbb{Z}/N$,
for $\gamma=-1$ by the cosets of the odd residues modulo $2N$ --- so $R^*$ is
typically far larger than the $N$ shifts, which is precisely the extra symmetry the
framework exploits. For $\gamma=0$ the algebra is local: a polynomial is a unit iff
its constant term is nonzero, so $|\Sigma|=q^{N-1}$, and an automorphism sends $x$ to
any nilpotent generator $y_1x+\dots+y_{N-1}x^{N-1}$ with $y_1\ne 0$, giving
$|\Aut(R)|=(q-1)q^{N-2}$.

\medskip

Table~\ref{tab:symmetry} contrasts the symmetry structures; the three quotient
families share one row, differing only in the constant $\gamma$.

\begin{table}[h]
  \centering
  \setlength{\tabcolsep}{10pt}
  \begin{tabular}{lll}
    \toprule
    problem                            & store $\Sigma$        & query $\Kappa$                                \\
    \midrule
    matrix-mult $\mmt{l}{m}{n}$        & $\GL_m$ (right mult.) & $\GL_l$ (left mult.), and transpose ($l=m=n$) \\
    full-poly-mult                     & $\PGL_2(\F_q)$        & $\Gal(\F_q/\F_p)$                             \\
    poly-mult $\F_q[x]/(x^N{-}\gamma)$ & $R^*/\F_q^*$          & $\Aut(R)\rtimes\Gal(\F_q/\F_p)$               \\
    \bottomrule
  \end{tabular}
  \caption{First-argument symmetry factorizations; the last row covers the cyclic
    ($\gamma=1$), truncated ($\gamma=0$), and negacyclic ($\gamma=-1$) products. In every family the
    query elements are indexed so that the $i$-th element is $i$, which is why a
    certificate can name a symmetry by a single integer.}
  \label{tab:symmetry}
\end{table}

\section{Orbit enumeration}\label{sec:orbits}

The first stage enumerates exactly one canonical representative per orbit of
constraint subspaces, for every dimension $d=0,1,\dots,N_A$. Subspaces are built
up one constraint at a time and deduplicated with the same meet-in-the-middle trick.

Representatives are generated dimension by dimension, in lexicographic
order of their RREF. A dimension-$d$ subspace is formed by
appending to a dimension-$(d-1)$ representative one new constraint whose leading
pivot lies strictly above the existing pivots; candidates with a nonzero
entry at an existing pivot are skipped because Gauss--Jordan elimination would
reduce them against the current basis. Because candidates are visited in global
lexicographic order, the first time an orbit is seen its representative is
automatically the lexicographically least, so no separate minimality test is
needed.

Deduplication uses a hash set of RREF keys.
For each kept representative $c$ we insert the RREF of every store image
$\sigma(c)$, $\sigma\in\Sigma$; for each new candidate $q$ we probe with
the RREF of every query image $\kappa(q)$, $\kappa\in\Kappa$, and
discard $q$ if any probe hits, since a hit means $q\equiv (\kappa^{-1}\sigma)(c)$
lies in $c$'s orbit. Algorithm~\ref{alg:enum}, in Appendix~\ref{app:enum}, states the procedure.

\section{Rank lower-bound techniques}\label{sec:techniques}

The dynamic program sweeps subspace
dimension from $N_A$ down to $0$. At each dimension it processes all orbits of that
dimension in parallel, then inserts their freshly computed bounds into the orbit
map so that the next (smaller) dimension can consult them. This ordering is sound
because the two techniques that look anything up --- degenerate reduction and
backtracking --- only ever query \emph{strictly larger} dimensions, which have
already been settled. Each orbit tries four techniques in this order: flatten,
degenerate reduction, forced product, and backtracking; later techniques run only
if the earlier ones have not already met the goal.

\paragraph{Flatten.} Compute the three flattening matrix ranks of the constrained
tensor and take the maximum. This is
the base case and needs no other orbit.

\paragraph{Degenerate reduction.} Add one more independent constraint to the
current subspace, producing a strictly smaller subspace whose orbit was processed
at the previous (larger) dimension; its bound is a valid bound for the current
orbit. We scan candidate constraints, look up each enlarged subspace in the orbit
map, and keep the best. The proof
records the extra constraint together with the witness pair from
Equation~\eqref{eq:witness}.

\paragraph{Forced product.} This automates the Hopcroft--Kerr substitution~\cite{hopcroft1971minimizing}.
Grouping the tensor along the $C$ factor into $t$
slices, the slices whose $A\otimes B$ part has matrix rank $1$ are single products;
we greedily take a maximal linearly independent set of $s$ of them and assume an
optimal decomposition computes them literally. Stripping these $s$ terms leaves
$s(t-s)$ unknown field coefficients in the residual; we enumerate all $q^{s(t-s)}$
assignments, flatten each residual, and take the minimum, so the bound is
$s+\min$. We repeat for all three cyclic orientations of the tensor and keep the
best. Because the minimum
must range over \emph{every} assignment for soundness, and the coefficients live in
$\F_q$, for $q>2$ each coefficient runs over all $q$ values; we skip the technique
when $q^{s(t-s)}$ is too large to enumerate. The proof records only which of the three
orientations won.

\paragraph{Substitution with backtracking.} This is the workhorse. Fixing an
optimal decomposition of length $r$, we repeatedly pick a linear form in the
$A$-variables, set it to zero, drop the terms it kills, and recurse on the
resulting more-constrained orbit; if a sequence of substitutions that removes
$\delta$ terms lands in an orbit of bound $b$ with $\delta+b\ge$ target, then
$r\ge$ target. The search is a DFS over canonical $A$-components --- the nonzero
linear forms supported away from the current pivots, one per coset.
The DFS depth is capped at the known lower bound given by the previous techniques.
Several optimizations make it scale: the DFS branches are run in parallel;
a thread-local cache memoizes the bound of each enlarged subspace, halving
itself stochastically when it grows too large; the target bound is raised one unit
at a time so that easy branches prune early; and a step limit aborts large
branches. Each branch that reaches the target contributes one record to the proof:
the depth, the subset mask of substituted forms, and the witness pair identifying
the orbit it reduced to.

Appendix~\ref{sec:toy} works all four techniques on the toy example
$\Rk{\F_2}{\mmt{2}{2}{2}}\ge 7$.

\section{Proof certificates and verification}\label{sec:certificates}

Every bound is emitted as a certificate that a verifier rechecks.
The verifier is simpler than the search program, and it runs faster and uses less memory.

For each orbit, the certificate records
its canonical constraint subspace, the certified lower bound, and a proof that is
exactly one of the four techniques.
The bound for the unconstrained orbit is ``the'' bound for the problem.
Each proof stores only what the
verifier cannot cheaply recompute:
\begin{itemize}
  \item \textbf{Flatten:} nothing; the verifier recomputes the flattening rank.
  \item \textbf{Forced product:} the winning cyclic orientation, an integer in
        $\{0,1,2\}$.
  \item \textbf{Degenerate:} the extra constraint, together with the witness pair
        $(\kappa,\sigma)$ of Equation~\eqref{eq:witness} that identifies the enlarged
        subspace's orbit.
  \item \textbf{Backtracking:} for each DFS leaf that reached the target, its depth,
        the subset of substituted forms, and the witness pair of the orbit it reduced
        to.
\end{itemize}

The verifier replays the same sweep from dimension $N_A$ down to $0$, so each orbit
is checked only against strictly larger orbits already confirmed. It rechecks
flatten and forced product by recomputation, and degenerate and backtracking by
reconstructing the canonical form of every claimed reduction from its witness via
Equation~\eqref{eq:witness} and confirming the looked-up bound supports the claim. A
witness that misses a known orbit, a recomputed bound that falls short, or a
miscounted trace aborts verification. Soundness follows by induction on decreasing
dimension: at dimension $N_A$ only flatten and forced product apply, and every
smaller-dimension step reduces to an orbit verified earlier.

As a global sanity check, every lower bound our framework proves is consistent with the
prior literature: in no case (Tables~\ref{tab:matrix}, \ref{tab:full}, \ref{tab:cyc},
\ref{tab:trunc}, and~\ref{tab:nega}) does a proved bound exceed the best previously known
upper bound, and in many small cases it coincides with the exact rank that is already
known. A soundness bug that inflated a bound would, across the dozens of tabulated cases,
is likely to produce a lower bound exceeding the known upper bound or
much smaller than the known lower bound.
\section{Conclusion and open problems}\label{sec:conclusion}

We have given a single automated framework for bilinear-complexity lower bounds
over finite fields, parameterized only by a tensor and a symmetry
group, and used it to improve the long-standing bound on $\mmt{3}{3}{3}$ over
$\F_2$ to $20$, to settle three more matrix formats, and to obtain eighteen new
lower bounds for polynomial multiplication --- across the
full, cyclic, truncated, and negacyclic products.
Every bound comes with a certificate that an auditable verifier
rechecks, typically faster than the search.

The source code for the search and the verifier, together with all certificates, is
available at \\
\url{https://github.com/wcgbg/tensor-rank-lower-bound}.

Several questions remain. The exact rank of $\mmt{3}{3}{3}$ over $\F_2$ is still
open in $[20,23]$. Scaling to $\mmt{4}{4}{4}$ is blocked by the sheer number of
constraint orbits and will need stronger symmetry reduction. For polynomial
multiplication, the full, cyclic, truncated, and negacyclic bounds should extend to
larger $N$ with more compute. A natural next family is the multiplication in the
extension-field $\F_{q^n}$ over $\F_q$,
computing $f\cdot g \bmod p(x)$ for an irreducible $p$ of degree $n$ over $\F_q$.
Finally, the framework applies verbatim to other finite fields; finding genuinely
new bounds over fields beyond $\F_2$ and $\F_3$ is left for future work.

\section*{Acknowledgements}
Thanks to Youming Qiao for helpful discussions on the tensor isomorphism problem.

\bibliography{refs}

\appendix

\section{Computational results}\label{app:results}

Table~\ref{tab:results} reports the search and verification details for each new bound.
The bounds themselves and their comparison to prior work appear in
Tables~\ref{tab:matrix}, \ref{tab:full}, \ref{tab:cyc}, \ref{tab:trunc},
and~\ref{tab:nega}.

\begin{table}[h]
  \centering
  \setlength{\tabcolsep}{4pt}
  \begin{tabular}{lllcllc}
    \toprule
    result                                    & machine$^\dagger$        & \shortstack{search\\time} & \shortstack{search\\cost} & \shortstack{cert.\\size} & \shortstack{verify\\time} & \shortstack{verify\\cost} \\
    \midrule
    $\Rk{\F_2}{\mmt{2}{3}{4}}\ge 19^\ddagger$ & MacBook Air              & $2$ hrs                   & --                        & --                       & $2$ hrs                   & --                        \\
    $\Rk{\F_2}{\mmt{3}{3}{3}}\ge 20$          & MacBook Air              & $40$ min                  & --                        & $32$ MiB                 & $3$ sec                   & --                        \\
    $\Rk{\F_2}{\mmt{3}{3}{4}}\ge 25$          & c8g.\{16,48\}xlarge$^\S$ & $4.6$ days                & $52$ USD                  & $600$ MiB                & $70$ min                  & $1$ USD                   \\
    $\Rk{\F_2}{\mmt{3}{4}{4}}\ge 29$          & c8g.48xlarge             & $4.9$ hrs                 & $5$ USD                   & $500$ MiB                & $3.4$ hrs                 & $3$ USD                   \\
    \midrule
    $\Rk{\F_2}{\Fullt{6}}\ge 16$              & c8g.48xlarge             & $11$ min                  & --                        & --                       & $10$ min                  & --                        \\
    $\Rk{\F_2}{\Fullt{7}}\ge 19$              & c8g.48xlarge             & $3$ hrs                   & $3$ USD                   & $9$ MiB                  & $2.6$ hrs                 & $3$ USD                   \\
    $\Rk{\F_2}{\Fullt{8}}\ge 21$              & c8g.48xlarge             & $6$ hrs                   & $6$ USD                   & $224$ MiB                & $3$ hrs                   & $3$ USD                   \\
    $\Rk{\F_3}{\Fullt{6}}\ge 14$              & c8g.48xlarge             & $35$ min                  & --                        & $6$ MiB                  & $31$ min                  & --                        \\
    \midrule
    $\Rk{\F_2}{\Cyct{7}}\ge 13$               & c8g.48xlarge             & --                        & --                        & --                       & --                        & --                        \\
    $\Rk{\F_2}{\Cyct{8}}\ge 19$               & c8g.48xlarge             & --                        & --                        & --                       & --                        & --                        \\
    $\Rk{\F_2}{\Cyct{10}}\ge 22$              & c8g.48xlarge             & --                        & --                        & $8$ MiB                  & --                        & --                        \\
    $\Rk{\F_3}{\Cyct{9}}\ge 19$               & c8g.48xlarge             & $3$ min                   & --                        & $3$ MiB                  & --                        & --                        \\
    \midrule
    $\Rk{\F_2}{\Trunct{6}}\ge 13$             & c8g.48xlarge             & --                        & --                        & --                       & --                        & --                        \\
    $\Rk{\F_2}{\Trunct{7}}\ge 16$             & c8g.48xlarge             & --                        & --                        & --                       & --                        & --                        \\
    $\Rk{\F_2}{\Trunct{8}}\ge 19$             & c8g.48xlarge             & --                        & --                        & --                       & --                        & --                        \\
    $\Rk{\F_2}{\Trunct{9}}\ge 21$             & c8g.48xlarge             & --                        & --                        & --                       & --                        & --                        \\
    $\Rk{\F_2}{\Trunct{10}}\ge 24$            & c8g.48xlarge             & $1$ min                   & --                        & $3$ MiB                  & --                        & --                        \\
    $\Rk{\F_2}{\Trunct{11}}\ge 27$            & c8g.48xlarge             & $7$ min                   & --                        & $22$ MiB                 & --                        & --                        \\
    $\Rk{\F_3}{\Trunct{8}}\ge 17$             & c8g.48xlarge             & --                        & --                        & --                       & --                        & --                        \\
    $\Rk{\F_3}{\Trunct{9}}\ge 19$             & c8g.48xlarge             & $2$ min                   & --                        & $2$ MiB                  & --                        & --                        \\
    $\Rk{\F_3}{\Trunct{10}}\ge 21$            & c8g.48xlarge             & $3.4$ hrs                 & $3$ USD                   & $20$ MiB                 & --                        & --                        \\
    \midrule
    $\Rk{\F_3}{\Negat{9}}\ge 19$              & c8g.48xlarge             & $3$ min                   & --                        & $2$ MiB                  & --                        & --                        \\
    \bottomrule
  \end{tabular}

  \smallskip
  {\footnotesize\raggedright
    -- Very small number.\\
    $^\dagger$ Machines: Apple MacBook Air M4 with 16 GB RAM; AWS c8g.16xlarge Spot
    Instance, Graviton4 processors, 64 vCPUs, 128 GiB memory; AWS c8g.48xlarge Spot
    Instance, Graviton4 processors, 192 vCPUs, 384 GiB memory.\\
    $^\ddagger$ Run on the equivalent $\mmt{3}{2}{4}$ format. The proof is
    forced-product-heavy, so CPU verification takes about as long as the search.
    The verification can be accelerated by an NVIDIA L20 GPU to approximately 1 minute.\\
    $^\S$ Search on c8g.16xlarge. Verification on c8g.48xlarge.\\
    \par
  }
  \caption{Search and verification details for each new lower bound.}
  \label{tab:results}
\end{table}

\section{Orbit-enumeration algorithm}\label{app:enum}

Algorithm~\ref{alg:enum} spells out the orbit enumerator of
Section~\ref{sec:orbits}: each dimension-$d$ representative is built by extending
a dimension-$(d-1)$ one with a new highest-pivot constraint, and the inner
duplicate test is the meet-in-the-middle probe of that section. The
$\textit{visited}$ set is cleared at the start of each dimension, since the
subspace dimension is an invariant of the group action.

\begin{algorithm}[h]
  \caption{Enumerate constraint-subspace orbits}
  \label{alg:enum}
  \begin{algorithmic}[1]
    \State $\textit{reps}[0]\gets[\varnothing]$
    \For{$d=1$ \textbf{to} $N_A$}
    \State $\textit{visited}\gets\varnothing$
    \For{each representative $c\in \textit{reps}[d-1]$ (in lex order)}
    \For{each candidate constraint $r$ above all pivots of $c$ (in lex order)}
    \State $q \gets c \cup \{r\}$
    \If{\textbf{not} \Call{Seen}{$q$}}
    \State append $q$ to $\textit{reps}[d]$
    \For{$\sigma\in\Sigma$}\ \ insert $\RREF(\sigma(q))$ into \textit{visited}\EndFor
    \EndIf
    \EndFor
    \EndFor
    \EndFor
    \State \Return \textit{reps}
    \Statex
    \Function{Seen}{$q$}
    \For{$\kappa\in\Kappa$}
    \If{$\RREF(\kappa(q))\in\textit{visited}$}\ \Return \textbf{true}\EndIf
    \EndFor
    \State \Return \textbf{false}
    \EndFunction
  \end{algorithmic}
\end{algorithm}

\subsection*{Orbit counts}\label{app:counts}

Table~\ref{tab:orbits} reports the number of constraint-subspace orbits Algorithm~\ref{alg:enum} finds
for each matrix multiplication format over $\F_2$ and $\F_3$. The count depends on whether the transpose
symmetry is in force, which happens exactly when the format is square ($l=m=n$);
the polynomial families are swept by the same procedure.

\begin{table}[h]
  \centering
  \setlength{\tabcolsep}{12pt}
  \begin{tabular}{llrr}
    \toprule
    format $l\times m$ & transpose? & \# orbits over $\F_2$ & \# orbits over $\F_3$          \\
    \midrule
    $2\times 2$        & yes        & $10$                  & $10$                           \\
    $2\times 2$        & no         & $11$                  & $11$                           \\
    $2\times 3$        & no         & $31$                  & $31$                           \\
    $2\times 4$        & no         & $86$                  & $91$                           \\
    $3\times 3$        & yes        & $496$                 & $736$                          \\
    $3\times 3$        & no         & $710$                 & $1046$                         \\
    $3\times 4$        & no         & $158{,}426$           & ${>}1.6\times10^{6}\,^\dagger$ \\
    $4\times 4$        & yes        & ${>}1.6\times10^{11}$ & ${>}8.8\times10^{15}$          \\
    \bottomrule
  \end{tabular}

  \smallskip
  {\footnotesize\raggedright
    $^\dagger$ The exact value is computable but unlikely to improve any matrix
    bound.\par}
  \caption{Number of constraint-subspace orbits for matrix multiplication,
    $A=\F_q^{l\times m}$. ``transpose?'' marks whether the order-two transpose
    symmetry is included; it applies only to square formats $l=m=n$. The ${>}$
    entries are the counting estimate of Equation~\eqref{eq:orbit-count}.}
  \label{tab:orbits}
\end{table}

Where the exact count is out of reach, an orbit-counting (Burnside) estimate gives
its order of magnitude. The number of orbits is at least the number of constraint
subspaces divided by the group order; approximating the number of $d$-dimensional
subspaces of $\F_q^{lm}$ and summing over $d$ yields
\begin{equation}\label{eq:orbit-count}
  \sum_{d=0}^{lm}
  \frac{q^{lm\,d}}{\,|\GL_l(\F_q)\times\GL_m(\F_q)\times\GL_d(\F_q)|\cdot|T|\,},
\end{equation}
where $T=C_2$ when the transpose symmetry is present ($l=m=n$) and $T=C_1$
otherwise, and $|\GL_n(\F_q)|=\prod_{k=0}^{n-1}(q^n-q^k)$. This produces the
parenthesized entries of Table~\ref{tab:orbits}. Both the estimate and the
enumeration can be halved by duality: orthogonal complement matches each orbit of
$d$-dimensional subspaces with one of dimension $lm-d$ in equal number, so only the
dimensions $d\le\lfloor lm/2\rfloor$ need be enumerated.

\subsection*{More invariants}\label{app:invariants}

Clearing $\textit{visited}$ per dimension exploits a single $G$-invariant, the
subspace dimension: two subspaces in the same orbit always have the same
dimension, so the set never needs to hold subspaces of different dimensions
simultaneously. The same idea applies to \emph{any} $G$-invariant. Fixing a value
(or a tuple of values) of one or more invariants, keeping in $\textit{visited}$
only the representatives realizing that value, and sweeping the values one at a
time replaces a single large hash set with several small ones and never merges
orbits across values, since equivalent subspaces agree on every invariant. We do
not currently exploit invariants beyond the dimension, as enumeration memory is
not the bottleneck of our results; the invariants of $G$ available for each family
include the following.

\paragraph{Matrix multiplication ($A=\F_q^{l\times m}$, $G=(\GL_l(\F_q)\times\GL_m(\F_q))\rtimes C_2$).}
\begin{itemize}
  \item \textbf{Rank distribution:} $\rho_i(S)=|\{M\in S:\rank M=i\}|$ for
        $i=0,\dots,\min(l,m)$. The sandwich action $M\mapsto PMQ^{-1}$ and the
        transpose preserve the rank of each $M$, so $\rho$ is constant on orbits.
  \item \textbf{Point profiles:} let $B_0,\dots,B_{d-1}$ be a basis of $S$. The
        left profile is
        $\lambda_i(S)=|\{x\in\F_q^{l}:\rank[x^\top B_0,\dots,x^\top B_{d-1}]=i\}|$ for
        $i=0,\dots,m$, and the right profile is
        $\mu_i(S)=|\{y\in\F_q^{m}:\rank[B_0y,\dots,B_{d-1}y]=i\}|$ for $i=0,\dots,l$.
        Left multiplication preserves $\lambda$, right multiplication preserves $\mu$,
        and a change of basis of $S$ preserves both; when $l=m=n$ the transpose swaps
        them, so the unordered pair $\{\lambda,\mu\}$ is invariant.
\end{itemize}

\paragraph{Full polynomial multiplication ($A=$ degree-$(N-1)$ binary forms, $G_{\mathrm{full}}=\mathrm{P\Gamma L}_2(\F_q)$).}
A nonzero form is an effective degree-$(N-1)$ divisor on $\mathbb{P}^1$, and
$\mathrm{P\Gamma L}_2$ permutes the closed points of $\mathbb{P}^1$ preserving
their degree.
\begin{itemize}
  \item \textbf{Factorization-type distribution} (the analog of the rank
        distribution): for a nonzero $f\in S$, let $\tau(f)$ be the multiset of
        pairs $(\deg P,\operatorname{ord}_P f)$ over the closed points $P$ in the
        support of $\operatorname{div}(f)$ --- equivalently, the multiset of
        (degree, multiplicity) of the $\F_q$-irreducible factors of $f$. Every
        $g\in\mathrm{P\Gamma L}_2$ sends $f$ to a form with the same type, so
        $\rho_\tau(S)=|\{f\in S:\tau(f)=\tau\}|$ is invariant.
  \item \textbf{Point profile} (the analog of the point profiles): for a closed
        point $P$ and $k\ge1$, let $S_{P,\ge k}=\{f\in S:\operatorname{ord}_P f\ge
          k\}$; the local vanishing profile is the codimension sequence
        $v_P(S)=(\dim S-\dim S_{P,\ge k})_{k\ge1}$. Since $\mathrm{P\Gamma L}_2$
        permutes the degree-$\delta$ closed points, for each $\delta$ the multiset
        $\{v_P(S):\deg P=\delta\}$ is invariant; its $k=1$ term is the rank of the
        evaluation map from $S$ to the fiber at $P$. The projective action fuses
        $0$, $\infty$, and every other point into one $\mathbb{P}^1$, so here there is
        a single point profile rather than the left/right pair of matrix
        multiplication.
\end{itemize}

\paragraph{Cyclic polynomial multiplication ($A=R=\F_q[x]/(x^N-1)$, $G_{\mathrm{cyc}}=(R^*\rtimes\Aut(R)\rtimes\Gal(\F_q/\F_p))/\F_q^*$).}
By the Chinese remainder theorem $R\cong\prod_{i=1}^{t}\F_{q^{d_i}}$ over the
$q$-cyclotomic cosets of $\mathbb{Z}/N$, with projections
$\pi_i\colon R\to\F_{q^{d_i}}$ (evaluation at the $N$-th roots of unity).
$G_{\mathrm{cyc}}$ scales each $\pi_i$ by a unit, permutes the equal-degree
factors, and twists them by field automorphisms, so it preserves the degree of
each factor and the zero pattern of $a\in R$.
\begin{itemize}
  \item \textbf{Support-type distribution} (the analog of the rank
        distribution): for $a\in S$, let
        $\operatorname{supp}(a)=\{d_i:\pi_i(a)\ne0\}$ be the multiset of degrees of
        the factors on which $a$ is supported. Then
        $\rho_\sigma(S)=|\{a\in S:\operatorname{supp}(a)=\sigma\}|$ is invariant.
  \item \textbf{Projection-dimension profile} (the cyclic point profile): for
        each factor $\dim_{\F_q}\pi_i(S)\in\{0,\dots,d_i\}$ (a unit scales
        $\pi_i(S)$ within $\F_{q^{d_i}}$, preserving its $\F_q$-dimension); for each
        $\delta$ the multiset $\{\dim_{\F_q}\pi_i(S):d_i=\delta\}$ is invariant. These
        factors are exactly the evaluation points at the $N$-th roots of unity.
\end{itemize}

All three point profiles are the same construction --- evaluate the subspace at
the points the symmetry group permutes and record the resulting rank distribution
--- and differ only in the point set: the left/right vector spaces for matrix
multiplication, the projective line $\mathbb{P}^1$ for full polynomial
multiplication, and the $N$-th roots of unity for cyclic convolution.

\section{A worked example: $\Rk{\F_2}{\mmt{2}{2}{2}}\ge 7$}\label{sec:toy}

We illustrate the dynamic program on $\mmt{2}{2}{2}$ over $\F_2$, where $A$ is the
space of $2\times 2$ matrices and the symmetry group is
$G=(\GL_2(\F_2)\times\GL_2(\F_2))\rtimes C_2$ acting by $X\mapsto PXQ^{-1}$ and by
transposition. There are ten orbits of constraint subspaces. The program processes
them from the most constrained to the unconstrained tensor, assigning each a lower
bound by whichever of the four techniques of Section~\ref{sec:techniques} succeeds
first. We write $a_{ij}$ for the free entries of $A$ that survive the constraints,
and the constrained tensor is $\sum_{i,j,k} a_{ij}\otimes b_{jk}\otimes c_{ki}$
restricted to them.

\subsection*{Orbit 0: $\{a_{00}=0,\ a_{01}=0,\ a_{10}=0,\ a_{11}=0\}$}
$A=
  \begin{pmatrix}0 & 0 \\0&0
  \end{pmatrix}$. The constrained tensor is $T_0=0$, so
$\R(T_0)=0$.

\subsection*{Orbit 1: $\{a_{00}=0,\ a_{01}=0,\ a_{10}=0\}$}
$A=
  \begin{pmatrix}0 & 0 \\0&a_{11}
  \end{pmatrix}$, with tensor
$T_1=a_{11}\otimes b_{10}\otimes c_{01}+a_{11}\otimes b_{11}\otimes c_{11}$.
Flattening the $A$ and $B$ factors into one yields a matrix of rank $2$, so
$\R(T_1)\ge 2$.

\subsection*{Orbit 2: $\{a_{00}=0,\ a_{01}+a_{10}=0,\ a_{11}=0\}$}
$A=
  \begin{pmatrix}0 & a_{01} \\a_{01}&0
  \end{pmatrix}$, with tensor
$T_2=a_{01}\otimes b_{00}\otimes c_{01}+a_{01}\otimes b_{01}\otimes c_{11}
  +a_{01}\otimes b_{10}\otimes c_{00}+a_{01}\otimes b_{11}\otimes c_{10}$.
Flattening yields a matrix of rank $4$, so $\R(T_2)\ge 4$.

\subsection*{Orbit 3: $\{a_{00}=0,\ a_{01}=0\}$}
$A=
  \begin{pmatrix}0 & 0 \\a_{10}&a_{11}
  \end{pmatrix}$, with tensor
$T_3=a_{10}\otimes b_{00}\otimes c_{01}+a_{10}\otimes b_{01}\otimes c_{11}
  +a_{11}\otimes b_{10}\otimes c_{01}+a_{11}\otimes b_{11}\otimes c_{11}$.
Flattening yields rank $4$, so $\R(T_3)\ge 4$.

\subsection*{Orbit 4: $\{a_{00}=0,\ a_{01}+a_{10}=0\}$}
$A=
  \begin{pmatrix}0 & a_{01} \\a_{01}&a_{11}
  \end{pmatrix}$, with tensor
\[
  T_4=a_{01}\otimes b_{00}\otimes c_{01}+a_{01}\otimes b_{01}\otimes c_{11}
  +a_{01}\otimes b_{10}\otimes c_{00}+a_{01}\otimes b_{11}\otimes c_{10}
  +a_{11}\otimes b_{10}\otimes c_{01}+a_{11}\otimes b_{11}\otimes c_{11}.
\]
Flattening yields only $4$, so we invoke the forced-product technique
(Section~\ref{sec:techniques}), which automates a substitution of Hopcroft and
Kerr~\cite{hopcroft1971minimizing}. The technique rests on the following lemma.

\begin{lemma}[Hopcroft and Kerr~{\cite[Lemma~2]{hopcroft1971minimizing}}]\label{lem:forced-product}
  Let $f_0,\dots,f_{s-1},f_{s},\dots,f_{t-1}$ be bilinear forms of which $f_0,\dots,f_{s-1}$ are
  linearly independent and each is a single product (a form of rank $1$). If the
  whole set can be computed with $r$ multiplications, then it can be computed with
  $r$ multiplications $s$ of which are exactly $f_0,\dots,f_{s-1}$.
\end{lemma}

\begin{proof}[Proof sketch]
  Take any bilinear algorithm with $r$ products $m_1,\dots,m_r$ computing the
  $f_i$. Each $f_j$ is a fixed $\F_q$-linear combination of the $m_\lambda$. Since
  $f_0$ has rank $1$ and is independent of $f_1,\dots,f_{s-1}$, some product
  appearing in $f_0$ can be replaced by $f_0$ itself without shrinking the span of
  the products, so $f_0$ is then computed by one multiplication. Repeating for
  $f_1,\dots,f_{s-1}$ --- whose independence guarantees a fresh product to trade at
  each step --- yields an $r$-product algorithm that computes $f_0,\dots,f_{s-1}$
  literally.
\end{proof}

Fix an optimal decomposition $T_4=\sum_{\lambda} u_\lambda\otimes v_\lambda\otimes
  w_\lambda$ with $r=\R(T_4)$ terms. Grouping $T_4$ along the $C$ factor exposes its
slices as forms in $A\otimes B$: the slice at $c_{00}$ is the single product
$a_{01}\otimes b_{10}$ and the slice at $c_{10}$ is the single product
$a_{01}\otimes b_{11}$, and these two are linearly independent (so $s=2$). By
Lemma~\ref{lem:forced-product} we may assume the decomposition computes them
literally, say $u_0\otimes v_0=a_{01}\otimes b_{10}$ and
$u_1\otimes v_1=a_{01}\otimes b_{11}$.

Now strip these two terms from $T_4$. The lemma pins down their $A\otimes B$ parts
but not their $C$-components $w_0,w_1$: each must match the slice it computes along
$c_{00}$ (resp.\ $c_{10}$), but its coordinates along the remaining basis vectors
$c_{01},c_{11}$ are unconstrained. Over $\F_2$ this leaves four unknown bits
$\mu_0,\dots,\mu_3$, and the residual is
\[
  T_4'=T_4+a_{01}\otimes b_{10}\otimes(c_{00}+\mu_0 c_{01}+\mu_1 c_{11})
  +a_{01}\otimes b_{11}\otimes(c_{10}+\mu_2 c_{01}+\mu_3 c_{11}),
  \qquad \mu_i\in\{0,1\}.
\]
Enumerating all $2^4=16$ assignments of $\mu$ and flattening each gives $\R(T_4')\ge 4$ in
every case, so $\R(T_4)\ge 4+2=6$.

\subsection*{Orbit 5: $\{a_{00}=0,\ a_{11}=0\}$}
$A=
  \begin{pmatrix}0 & a_{01} \\a_{10}&0
  \end{pmatrix}$, with tensor
$T_5=a_{01}\otimes b_{10}\otimes c_{00}+a_{01}\otimes b_{11}\otimes c_{10}
  +a_{10}\otimes b_{00}\otimes c_{01}+a_{10}\otimes b_{01}\otimes c_{11}$.
Flattening yields rank $4$, so $\R(T_5)\ge 4$.

\subsection*{Orbit 6: $\{a_{01}+a_{10}=0,\ a_{00}+a_{01}+a_{11}=0\}$}
$A=
  \begin{pmatrix}a_{00} & a_{01} \\a_{01}&a_{00}+a_{01}
  \end{pmatrix}$. Flattening gives
only $4$. Here, we try to get a better lower bound using substitution with backtracking (Section~\ref{sec:techniques}).
Fix an optimal decomposition $T_6=\sum_\lambda u_\lambda\otimes v_\lambda\otimes
  w_\lambda$. The free entries span $\{a_{00},a_{01}\}$, whose only nonzero linear
forms over $\F_2$ are $a_{00}$, $a_{01}$, and $a_{00}+a_{01}$, so every nonzero
$u_\lambda$ is one of these three. Since the decomposition has at least $4$ terms,
by pigeonhole one of the three forms occurs as some $u_\lambda$ at least twice, and
setting it to zero kills at least two terms:
\begin{itemize}
  \item $a_{00}=0$ leaves $
          \begin{pmatrix}0 & a_{01} \\a_{01}&a_{01}
          \end{pmatrix}$,
        Orbit 2 after adding row $0$ to row $1$;
  \item $a_{01}=0$ leaves $
          \begin{pmatrix}a_{00} & 0 \\0&a_{00}
          \end{pmatrix}$,
        Orbit 2 after swapping the rows;
  \item $a_{00}+a_{01}=0$ leaves $
          \begin{pmatrix}a_{00} & a_{00} \\a_{00}&0
          \end{pmatrix}$,
        Orbit 2 after adding row $1$ to row $0$.
\end{itemize}

This pigeonhole argument of the lower bound 6 can be automated by the backtracking below.
\begin{itemize}
  \item $a_{00}$: $6$ not met; go deeper.
        \begin{itemize}
          \item $a_{00},a_{00}$: set $a_{00}=0$, reaching Orbit 2; $\R(T_6)\ge\R(T_2)+2\ge 6$. Proved; backtrack.
          \item $a_{00},a_{01}$: $6$ not met; go deeper.
                \begin{itemize}
                  \item $a_{00},a_{01},a_{01}$: set $a_{01}=0$, reaching Orbit 2; $\R(T_6)\ge 6$. Proved; backtrack.
                  \item $a_{00},a_{01},a_{00}+a_{01}$: $6$ not met; go deeper.
                        \begin{itemize}
                          \item $a_{00},a_{01},a_{00}+a_{01},a_{00}+a_{01}$: set $a_{00}+a_{01}=0$, reaching Orbit 2; $\R(T_6)\ge 6$. Proved; backtrack.
                        \end{itemize}
                \end{itemize}
          \item $a_{00},a_{00}+a_{01}$: $6$ not met; go deeper.
                \begin{itemize}
                  \item $a_{00},a_{00}+a_{01},a_{00}+a_{01}$: set $a_{00}+a_{01}=0$, reaching Orbit 2; $\R(T_6)\ge 6$. Proved; backtrack.
                \end{itemize}
        \end{itemize}
  \item $a_{01}$: $6$ not met; go deeper.
        \begin{itemize}
          \item $a_{01},a_{01}$: set $a_{01}=0$, reaching Orbit 2; $\R(T_6)\ge 6$. Proved; backtrack.
          \item $a_{01},a_{00}+a_{01}$: $6$ not met; go deeper.
                \begin{itemize}
                  \item $a_{01},a_{00}+a_{01},a_{00}+a_{01}$: set $a_{00}+a_{01}=0$, reaching Orbit 2; $\R(T_6)\ge 6$. Proved; backtrack.
                \end{itemize}
        \end{itemize}
  \item $a_{00}+a_{01}$: $6$ not met; go deeper.
        \begin{itemize}
          \item $a_{00}+a_{01},a_{00}+a_{01}$: set $a_{00}+a_{01}=0$, reaching Orbit 2; $\R(T_6)\ge 6$. Proved; backtrack.
        \end{itemize}
\end{itemize}
Every branch reaches the target, so $\R(T_6)\ge 6$.

\subsection*{Orbit 7: $\{a_{00}=0\}$}
$A=
  \begin{pmatrix}0 & a_{01} \\a_{10}&a_{11}
  \end{pmatrix}$. Adding the constraint
$a_{01}+a_{10}=0$ lands in Orbit 4, so degenerate reduction
(Section~\ref{sec:techniques}) gives $\R(T_7)\ge\R(T_4)\ge 6$.

\subsection*{Orbit 8: $\{a_{01}+a_{10}=0\}$}
$A=
  \begin{pmatrix}a_{00} & a_{01} \\a_{01}&a_{11}
  \end{pmatrix}$. Adding the constraint
$a_{00}=0$ lands in Orbit 4, so $\R(T_8)\ge\R(T_4)\ge 6$.

\subsection*{Orbit 9: $\{\,\}$}
$A=
  \begin{pmatrix}a_{00} & a_{01} \\a_{10}&a_{11}
  \end{pmatrix}$ is the unconstrained
tensor $T_9=\sum_{i,j,k}a_{ij}\otimes b_{jk}\otimes c_{ki}=\mmt{2}{2}{2}$. We apply
substitution with backtracking. Fix an optimal decomposition
$T_9=\sum_\lambda u_\lambda\otimes v_\lambda\otimes w_\lambda$. If $a_{00}$ appears
as some $u_\lambda$, setting $a_{00}=0$ lands in Orbit 7 and gives
$\R(T_9)\ge\R(T_7)+1\ge 7$; by the sandwich symmetry the same holds when the form
is any of $a_{01}$, $a_{10}$, $a_{11}$, $a_{00}+a_{01}$, $a_{10}+a_{11}$,
$a_{00}+a_{10}$, $a_{01}+a_{11}$, or $a_{00}+a_{01}+a_{10}+a_{11}$. Otherwise the
form is one of $a_{01}+a_{10}$, $a_{00}+a_{11}$, $a_{00}+a_{01}+a_{10}$,
$a_{00}+a_{01}+a_{11}$, $a_{00}+a_{10}+a_{11}$, or $a_{01}+a_{10}+a_{11}$, and
setting it to zero lands in Orbit 8, giving $\R(T_9)\ge\R(T_8)+1\ge 7$. These
exhaust the nonzero forms, so $\Rk{\F_2}{\mmt{2}{2}{2}}=\R(T_9)\ge 7$.

\end{document}
